% ==============================================================================
% Plantilla Institucional de Trabajo Terminal — UPIIZ IPN
% Archivo: capitulos/protocolo/07-metodologia-trabajo.tex
% ==============================================================================
% ESTRUCTURA NORMATIVA DE PROTOCOLO: CAPÍTULO 7 - METODOLOGÍA DE TRABAJO
% ==============================================================================
% LÍMITE INSTITUCIONAL: Máximo 5 páginas.
%
% Requisitos normativos (Anexo 1):
% 1. Identificación y análisis del problema.
% 2. Listado preliminar de especificaciones de diseño.
% 3. Posibles soluciones.
% 4. Metodología de diseño mecatrónico aplicada al proyecto.
% 5. Etapa prevista de implementación.
% 6. Etapa prevista de pruebas.
%
% NOTA DE TIEMPO VERBAL:
% Redactar estrictamente en modo prospectivo o futuro (lo que se prevé realizar),
% no en pasado como si el diseño, implementación o pruebas ya hubieran ocurrido.
% ==============================================================================

\chapter{Metodología de trabajo}
\label{cap:proto-metodologia}

% [INSTRUCCIÓN PARA EL ALUMNO]:
% Desarrolle este capítulo en un máximo de 5 páginas describiendo la metodología que se seguirá:

\section{Identificación y análisis del problema}

[Presente la identificación detallada y el análisis riguroso de la problemática técnica a resolver, explicando las variables involucradas y los factores operativos que se considerarán].

\section{Listado preliminar de especificaciones de diseño}

[Enuncie el listado de especificaciones preliminares cuantitativas y cualitativas que deberá cumplir el diseño (rangos de operación, tolerancias, dimensiones, capacidades, consumos o requerimientos normativos)]:

\begin{itemize}
    \item \textbf{Especificación 1:} [Descripción de la especificación técnica preliminar y valor objetivo esperado].
    \item \textbf{Especificación 2:} [Descripción de la especificación técnica preliminar y valor objetivo esperado].
    \item \textbf{Especificación 3:} [Descripción de la especificación técnica preliminar y valor objetivo esperado].
\end{itemize}

\section{Posibles soluciones}

[Exponga las alternativas preliminares o posibles soluciones técnicas analizadas para resolver el problema y fundamente la selección de la opción que se proyecta desarrollar].

\section{Metodología de diseño mecatrónico aplicada al proyecto}

[Describa la metodología de diseño mecatrónico que se aplicará al desarrollo del proyecto, indicando cómo interactuarán las etapas de modelado, diseño mecánico, diseño electrónico, control y programación].

\section{Etapa prevista de implementación}

[Explique cómo se tiene prevista la etapa de implementación, manufactura, adquisición, programación y ensamble físico o lógico del sistema a lo largo del desarrollo].

\section{Etapa prevista de pruebas}

[Detalle el protocolo de pruebas experimentales que se prevé realizar para evaluar el desempeño y funcionamiento del sistema una vez implementado].
